\section{Dataset Format and Context Taxonomy}

A central failing of historical biomechanics datasets is a lack of rigorous contextual metadata. A raw $1\,\text{kHz}$ acceleration array is analytically worthless if the researcher does not know the mass of the barbell, the athlete's fatigue state, the calibration of the load, or the specific intent of the training block. 

To address this, our system implements a highly formalized, schema-versioned (Schema v4) JSON format that deeply integrates with the Brain Imaging Data Structure (BIDS) layout principles for multi-subject archiving.

\subsection{The Session JSON Object}
Every recorded session generates a root \texttt{session\_id.json} file. This file acts as the primary key and provenance chain for all raw CSV files in that folder. The structure is broken down into modular structs defined in the C++ \texttt{Config.h}.

\subsubsection{Subject Identity and Anthropometrics}
To comply with ethics protocols while maintaining longitudinal tracking, the system uses an anonymized UUID (`subject_uuid`). This UUID maps consistently across all sessions for a specific athlete. 
The \texttt{SubjectDaySnapshot} object captures the state of the athlete \textit{on that specific day}. This is vital, as biological organisms are non-stationary.
\begin{itemize}
    \item \texttt{body\_mass\_kg}: The athlete's weight that morning.
    \item \texttt{sleep\_hours\_last\_night} \& \texttt{sleep\_quality\_1to5}: Drives the noise floor of the athlete's force output.
    \item \texttt{caffeine\_mg}: Quantifies nervous system stimulation.
    \item \texttt{soreness\_1to10} \& \texttt{soreness\_locations}: Useful for correlating asymmetric bar paths with localized muscular trauma.
\end{itemize}

\subsubsection{Training Context}
Kinematics are entirely dependent on intent. A $100\,\text{kg}$ squat moves slowly if the athlete intends to move it slowly (e.g., a hypertrophy block) and rapidly if the intent is power development. The \texttt{TrainingContext} object records:
\begin{itemize}
    \item \texttt{block\_phase}: Accumulation, Intensification, Realization, Deload.
    \item \texttt{tempo\_prescription}: E.g., "2-0-1-0" denoting a 2-second eccentric, no pause, explosive concentric, and no pause.
    \item \texttt{prior\_session\_volume\_kg}: The total tonnage lifted in the previous session, proxying cumulative peripheral fatigue.
\end{itemize}

\subsubsection{Load Provenance and Gear}
Knowing the exact load on the barbell is critical for calculating Force ($F=ma$) and Power ($P=Fv$).
The \texttt{LoadProvenance} array breaks down the exact plates loaded on the bar (e.g., `[20, 10, 5, 2.5]`). It tracks the `collar_mass_kg_each`, and includes a boolean flag `plates_calibrated` indicating if the plates were weighed on a calibrated scale rather than relying on manufacturer stamping.

The \texttt{GearAndImplements} struct meticulously tracks supportive equipment:
\begin{itemize}
    \item \texttt{belt\_type}: Lever, Prong, Velcro. Belts massively increase intra-abdominal pressure, altering torso angle and bar trajectory.
    \item \texttt{shoe\_type}: Heeled lifters vs flats. Heeled shoes alter the ankle dorsiflexion angles, heavily modifying the kinematic chain.
    \item \texttt{knee\_wraps}: Wraps store elastic potential energy during the eccentric phase, creating an artificial velocity spike at the bottom of the concentric phase.
\end{itemize}

\subsubsection{Subjective Quality}
Post-set, the athlete provides subjective feedback. The \texttt{SubjectiveQuality} struct logs the Rating of Perceived Exertion (RPE on a 1-10 scale), the true Repetitions in Reserve (RIR), and free-text `technique_breakdown_notes` (e.g., "lower back rounded on rep 4"). This allows researchers to train Machine Learning models to correlate specific IMU signal signatures with form breakdown and RPE.

\subsection{Multi-Set Architecture}
Earlier versions of the schema assumed one file per set. Schema v4 introduces a nested \texttt{std::vector<SetInfo> sets} array. A single session directory now holds an entire workout. Each entry in the `sets` vector documents the `t_start_unified_s` and `t_end_unified_s`, the `barbell_weight_kg`, the `target_reps`, and the `completed_reps`. This unified timeline allows researchers to analyze inter-set recovery periods and the compounding degradation of kinematics over a full 60-minute training block.
